\documentclass{book}

%*****************
\input{structure} 
\begin{document}
\newcommand{\theimage}{fig22.jpg}
\pagestyle{empty}

% wstawienie numerów stron w ramki dynamiczne frameD-1a i frameD-1b
\begin{dynamiccontents*}{frameD-1a}
\begin{tikzpicture}
\draw(0,0) node [fill=cafe, minimum width=2cm, minimum height=2cm]{
{\sffamily\bfseries\Huge\color{white}\thepage}
};
\end{tikzpicture}
\end{dynamiccontents*}

\begin{dynamiccontents*}{frameD-1b}
\begin{tikzpicture}
\draw(0,0) node [fill=cafe, minimum width=2cm, minimum height=2cm]{
{\sffamily \bfseries\Huge\color{white}\thepage}
};
\end{tikzpicture}
\end{dynamiccontents*}


% wstawienie grafiki w ramkię statyczną frameS-1a
\begin{staticcontents*}{frameS-1a}
\includegraphics[width=1\textwidth]{\theimage}
\end{staticcontents*}

% wypełnienie tekstem ramki statycznej frameS-1b
\begin{staticcontents*}{frameS-1b}
\begin{tikzpicture}
\draw(0,0) node [fill=white, text width=13cm, inner sep=5mm, opacity=0.7]{
\large\sffamily
{\fontsize{70}{30}\selectfont {\color{black} RECANEWS}}\\


{\fontsize{30}{30}\selectfont {\color{black} Volumen 4}}\\
\begin{flushright}
{\fontsize{30}{30}\selectfont {\color{black} Abril 2015}}
\end{flushright}
};
\end{tikzpicture}
\end{staticcontents*}


% wlanie tekstu do wszystkich ramek typu flow



\newpage


\renewcommand\thesection{\arabic{section}}
\renewcommand\thesubsection{\arabic{subsection}}

 
\tableofcontents{}
%*****************

%*********************************************
\addcontentsline{toc}{section}{Noticias Astronomía y ciencias del espacio en Colombia}
\section*{Noticias Astronomía y ciencias del espacio en Colombia}
%*********************************************
\subsection{LARIM Pre-Inscripción}

Ya esta implementada la plataforma de registro como participante de nuestra LARIM 2016
\begin{center}
\url{http://larim.unal.edu.co}
\end{center}
\textbf{¡Es de suprema importancia motivar a nuestros colegas y estudiantes a que hagan su pago durante la inscripción temprana!}
\begin{itemize}
\item Consignación
\end{itemize}
Consignar a nombre del titular Academia Colombiana de Ciencias Exactas, Físicas y Naturales\\
Nit: 860.026.6351\\
Cuenta de Ahorros: 055000 99 00 29 2575\\
Banco: DAVIVIENDA\\
País: Colombia\\
Una vez realizada su consignación, ingrese a la plataforma de registro, diligencie las casillas y  suba un archivo adjunto que contenga la imagen, claramente legible, del recibo de la consignación. 
\begin{itemize}
\item Costos
\end{itemize}
\textbf{Antes del 1 de Mayo de 2015}\\
Estudiantes (carnet vigente Octubre-2016): 250.000 pesos*\\
Profesionales 400.000 pesos*\\
*Incluye sin costo adicional la participación en la cena del evento.


\subsection{Exorings on the Horizon}

Astrónomos del Centro Harvard-Smithsonian para Astrofísica y de la Universidad de Antioquia (Medellin-Colombia), han ideado un método novedoso para identificar anillos alrededor de planetas extrasolares (exoanillos).  El método es relativalemente simple y puede ser usado para analizar grandes bases de datos fotométricas y encontrar así una lista de candidatos para futuros estudios.

Encuentre el artículo completo en

\begin{center}
\url{http://goo.gl/PXmAlh}
\end{center}

\subsection{Falling Walls Lab Bucaramanga}

\textbf{Descripción:}\\
\\
La convocatoria a Falling Walls Lab Bucaramanga está abierta para académicos, emprendedores y profesionales de todas las disciplinas que quieran presentar sus trabajos de investigación, proyectos de innovación, modelos de negocios o ideas de alto impacto. Cada participante tiene la oportunidad de realizar una presentación de 3 minutos (en inglés), frente a un jurado calificador interdisciplinar compuesto por expertos de diversas áreas.\\
\\
Durante el transcurso del año se desarrollan Falling Walls Labs en todo el mundo, en países como Marruecos, Rusia, Sudáfrica, Finlandia, Italia, Croacia y Estados Unidos. Los ganadores de estos concursos internacionales, clasifican a la final, el 8 de noviembre en Berlín. Este evento se lleva a cabo de forma conjunta con The Falling Walls Conference, que se realiza el 9 de noviembre, evento en el que 20 conferencistas, expertos internacionales en diversas áreas, presentan sus investigaciones y proyectos, tratando de dar respuesta a la pregunta: ``Which are the next walls to fall?''.\\
\\
\textbf{Requisitos:}\\
\\
Pueden participar académicos, estudiantes de posgrado, investigadores, emprendedores y profesionales de todas las disciplinas, que sean menores de 35 años.\\
\\
\textbf{Premio:}\\
\\
El ganador recibirá tiquetes, hospedaje y alimentación para participar en la final de Falling Walls Lab que se realizará en Berlín, Alemania el próximo 8 de noviembre. Los ganadores de esta final participarán como expositores en The Falling Walls Conference al día siguiente.\\
\\
\textbf{Fecha de cierre:}\\
\\
La convocatoria estará abierta hasta el próximo 2 de Mayo de 2015. Los participantes seleccionados para participar en la final del 28 de Mayo serán anunciados durante la segunda semana de Mayo. Más información\\
\begin{center}
\url{http://goo.gl/2kzRn7}
\end{center}

\subsection{Territorialización del espacio ultraterrestre}
Este foro, nace de la necesidad de generar un espacio de discusión en torno a la apropiación de la órbita geoespacial colombiana y de la importancia de producir herramientas propias de captura de información espacial actualizada y de libre acceso. Los objetivos a alcanzar son: 

    Visibilizar la importancia de generar estrategias transdisciplinarias que conlleven al desarrollo de la tecnología espacial colombiana, para impulsar políticas que conduzcan a una soberanía tecnológica
    Analizar y problematizar la territorialización del espacio desde las políticas del Estado, el acceso a la información, y las implicaciones de su uso en política pública.
    Motivar a la comunidad académica y al público general a tomar conciencia sobre estos temas.

\begin{description}
\item[Programa]
\item[Día 1:]  8 de Marzo\\
\item[*]El Espacio Ultraterrestre y la necesidad de su territorialización
\item[*]Desarrollo de la Astronáutica a través de la historia
\item[*]Los satélites y la experiencia colombiana


\item[Día 2:] 19 de Marzo
\item[*]La Comisión Colombiana del Espacio: ¿qué ha pasado en nueve años?
\item[*]Derecho al espacio: Colombia y la legislación espacial
\item[*]Colombia: ¿Hacia la soberanía tecnológica de su espacio ultraterrestre?\\
\item[Lugar:]SALÓN OVAL ED. POSGRADOS DE CIENCIAS HUMANAS\\
UNIVERSIDAD NACIONAL DE COLOMBIA
\item[Hora:] DE 4:00 A 8:00 PM
\end{description}

Luego de cada sesión se realizará una actividad de observación en el Observatorio Astronómico de la Universidad, el cual cuenta con un cupo de 40 personas. Las inscripciones para esta actividad se realizaran en al inicio de cada jornada.\\

Más información e inscripción:
\begin{center}
\url{http://goo.gl/lDuYI8}
\end{center}

\subsection{Posible posición para un Astrofísico en la UNAL}

La Facultad de Ciencias de la Universidad Nacional  acaba de iniciar proceso de convocatoria para proveer cargos docentes.\\
\\
Toda la información está en:
\begin{center}
\url{http://goo.gl/edNJzK}
\end{center}
En uno de ellos, el perfil E17 correspondiente al Departamento de Química, requiere un cargo para docencia en el área de la espectroscopia. Puede ser un profesional en el área de   química o  física y el doctorado ha de ser equivalente en el área de desempeño con tesis que tenga que ver en desarrollos de espectroscopia.\\
\\
ver:
\begin{center}
\url{http://goo.gl/ET1GKp}
\end{center}
Perfectamente algún astrofísico que haya hecho su tesis de doctorado en espectroscopia astronómica y que sean físicos de formación, pueden presentarse para tal perfil.

\subsection{Posición Postdoc en la UIS}

Se busca un Postdoc para experimentos: MuonTelescope (MuTe) y proyecto LAGO (\url{http://lagoproject.org})\\

Deadline: 2015-05-31\\

Más Información
\begin{center}
\url{https://inspirehep.net/record/1348393}
\end{center}

\subsection{1st workshop on Current Challenges in Cosmology: Inflation and the Origin of CMB anomalies.\\ Mayo 18 - 22, 2015 Cali, Colombia}
The Workshop is intended to gather leading experts in theoretical Cosmology to share knowledge, discuss new theoretical developments in inflationary cosmology and ultimately elucidate the puzzle of CMB anomalies. 
\textbf{Main topics}
\begin{itemize}
\item Statistical anisotropy and anisotropic expansion
\item Primordial magnetic fields
\item Parity violation and polarization in the CMB
\item Effective field theory of inflation
\item Anisotropic non-gaussianity
\item CMB data analysis
\item Primordial Gravitational  Waves
\item Galileons and Horndeski theories
\end{itemize}
\par
\noindent Invited Speakers:\\
\\
Konstantinos Dimopoulos (Lancaster University),	David F. Mota (University of Oslo), Maresuke Shiriashi (Università degli Studi di Padova), Patrick Peter
(Institut d’Astrophysique de Paris), Azadeh Maleknejad (Institute for Research in Fundamental Sciences), Marco Peloso (University of Minnesota), Lorenzo Sorbo (University of Massachusetts), Thiago S. Pereira (Universidade Estadual de Londrina), Leonardo Senatore* (Stanford University), Mindaugas Karciauskas (Helsinki University), Martin Bucher (University Paris 7 Diderot).\\
\\
Más información

\begin{center}
\url{cosmology.univalle.edu.co/index.php/home}
\end{center}

\subsection{Publicación - Juan C. Ibañez}
Al candidato a Doctor  Mejía del Institut fuer theoretische Astrophysik - Heidelberg Universitaet, le aceptaron su primer artículo en la revista la revista Astronomy \& Astrophysics. \\
\\
Este artículo sintetiza la mayor parte de la tesis de maestría de Juan realizada en la Universidad de Bonn en donde estuvo trabajando en el grupo de astrofísica estelar. En este proyecto se intentó comprender el comportamiento de la inestabilidad de Tayler en zonas radiativas en el interior de estrellas masivas y su influencia en el transporte del momento angular y de elementos producto de la evolución estelar entre el núcleo y la superficie de estas estrellas.\\
\\
El artículo completo puede ser consultado en
\begin{center}
\url{http://arxiv.org/abs/1503.04450}
\end{center}

\subsection{Publicación - Santiago Vargas}
El doctor Santiago Vargas, profesor del Observatorio Astronómico Nacional le fue aceptado una publicación en la revista especializada Solar Physics. El doctor Vargas explica resumen su publicación:\\

A partir de imágenes de alta resolución del telescopio Solar Swedish Telescope y el satélite Hinode, se estudia la evolución de vortices de plasma en la fotosfera solar. Se analiza el efecto de los vortices sobre la dinámica de estructuras magnéticas a pequeña escala (puntos brillantes).\\

Este es el link de la publicación
\begin{center}
\url{http://goo.gl/mUtpG7}
\end{center}
\newpage
\subsection{Astropuerta - Abril}
\begin{itemize}
\item Principales eventos celestes de Abril 2015
\end{itemize}
Sábado 4 - Luna llena
Sábado 4 - Eclipse total de Luna visible en Australia, Japón, Alaska y en el Océano Pacífico
Miércoles 8 - Conjunción de la Luna y Saturno
Domingo 12 - Luna en cuarto menguante
Sábado 18 - Luna nueva
Martes 21 - Ocultación de la estrella Aldebarán por la Luna visible en Canadá, Islandia y Escandinavia
Jueves 23 - Conjunción de Mercurio y Marte
Sábado 25 - Luna en cuarto creciente
Domingo 26 - Ocultación del asteroide Juno por la Luna visible en el Sureste de Asia

\begin{itemize}
\item Principales efemérides históricas de Abril 2015
\end{itemize} 
    Jueves 2 - 1845: Primera fotografía del Sol
    Viernes 3- 1966: La sonda Lunik 3, primera nave en orbitar la Luna
    Domingo 12 - 1961: Yuri Gagarin, primer hombre en el espacio
    Domingo 12 - 1981: Lanzamiento del Columbia, primer Transbordador Espacial
    Martes 14 - 1629: Nace Christiaan Huygens, descubridor de la verdadera forma del anillo de Saturno
    Viernes 17 - 2007: Lanzamiento de Libertad 1, primer satélite colombiano.
    Viernes 17 - 2014: Se confirma el descubrimiento del exoplaneta Kepler-186f, el primero similar a la Tierra en tamaño, composición y distancia a la estrella
    Domingo 19 - 1971: La Unión Soviética lanza la primera estación espacial, la Salyut 1
    Jueves 23 - 1967: Accidente mortal del cosmonauta Vladimir Komarov en la nave Soyuz 1         
    Viernes 24 - 1970: China lanza su primer satélite artificial
    Sábado 25 - 1990: Lanzamiento del Telescopio Espacial Hubble
    Martes 28 - 2001: Dennis de Tito, a bordo de la Estación Espacial Internacional, primer turista en el espacio         
    \begin{flushright}
    Germán Puerta
    \end{flushright}
 \newpage
%*********************************************
\addcontentsline{toc}{section}{Escuelas}
       \section*{Escuelas}
%*********************************************

\subsection{9$^{a}$ Escuela de Verano en Astrofísica.  Deadline 15 de Abril }

La escuela de Verano es una escuela introductoria a la astrofísica observacional y teórica, dirigida a estudiantes de quinto semestre en adelante de licenciaturas en Física, Matemáticas y áreas afines. Todas las actividades se realizarán dentro de las instalaciones del Centro de Radioastronomía y Astrofísica (CRyA), UNAM Campus Morelia.\\
\\
Habrá cursos introductorios, charlas especializadas y talleres interactivos sobre diversos temas de la astrofísica moderna, tanto observacionales como teóricos.\\
\\
Los temas que se tratarán durante la escuela incluyen: Cosmología, Astrofísica observacional, Astrofísica de altas energías, Astrofísica extra-galáctica, Formación estelar, Medio interestelar, Radioastronomía y Planetas extrasolares. Más información:
\begin{center}
\url{www.crya.unam.mx/verano/}
\end{center}

\subsection{XXIV Verano en el Observatorio 2015}
\begin{description}
\item[Objetivo]
Proveer un conocimiento teórico-práctico de diversos aspectos de la astronomía actual, así como un acercamiento a la investigación científica.
\item[Información]
El curso esta dirigido a estudiantes de licenciatura en física o en ciencias afines que estén interesados en adquirir un conocimiento amplio de la astrofísica actual y de las técnicas observacionales mas usadas en el óptico e infrarrojo.
Los estudiantes seleccionados tendrán la oportunidad de participar con un investigador en el desarrollo de un proyecto, el cual se presentara al final del curso. ademas podrán trabajar algunas noches con los telescopios del observatorio astronómico nacional.
\item[Requisitos]
1) La solicitud se llena en linea a través de esta pagina en la sección de solicitud.\\
2) Se requieren dos cartas de recomendación por parte de profesores de la institución de la cual procede el estudiante; estas cartas de recomendación se llenan en linea en la sección de solicitudes (sub-sección carta de recomendación).\\
3) No haber asistido a este curso de verano en ediciones pasadas.
\item[Solicitudes] 1 de Febrero al 20 de Abril
\end{description}
\begin{center}
\url{http://bufadora.astrosen.unam.mx/~verano/}
\end{center}

%*********************************************
\addcontentsline{toc}{section}{Programas de posgrados}
       \section*{Programas de posgrados}
%*********************************************

\subsection{Becas de Doctorado Pre-doctorales en el IAC - España}
El instituto de Astrofísica de Canarias tiene 7 posiciones pre-doctorales para desarrollar proyectos en simulaciones de la Vía Láctea  y comparación con datos de APOGEE, y de futuros datos de Gaia.

\begin{itemize}
\item CINCO contratos pre-doctorales en el marco del Programa\\
    "Astrofísicos Residentes del IAC":\\
\url{http://www.iac.es/info.php?op1=26&id=505}
\item DOS contratos pre-doctorales en el marco del "Programa Internacional\\
    de Doctorado LaCaixa-Severo Ochoa":\\
\url{http://www.iac.es/info.php?op1=26&id=506}
\end{itemize} 
\subsection{Posiciones doctorales por SFB 956}
The scientific goal is to understand the Cosmic Cycle and the cycle of stellar formation from the interstellar matter through life and deaths of stars. The positions offer a profound education in one of the following fields: experimental astrophysics, theoretical analyses and modeling, laboratory astrophysics and the development of new dedicated instrumentation. Furthermore, the successful candidates will be supported by the SFB 956 Doctoral Program.\\

Applications are invited for a PhD position within the CRC 956, a Collaborative Research Centre between the member institutions I. Physikalisches Institut, Universität zu Köln (UzK); Argelander-Institut für Astronomie, Universität Bonn (UB); Max-Planck-Institut für Radioastronomie (MPIfR). Depending on the PhD position the successful applicants should have a background in astro-physics, programming experience, interest in scientific computing or knowledge in engineering. All candidates should be fluent in English, German language skills are welcome but not conditional.For detailed information on project specific requirements 
\begin{center}
\url{http://www.sfb956.de/node/}
\end{center}

\newpage
\subsection{Posiciones Doctorales en USA para minorías}

Fellowships designed for underrepresented minorities and other relevant resources:
\begin{itemize}
\item Recruitment Fellowships: \url{http://goo.gl/88YgRz}
\item Recruitment Visits: \url{http://goo.gl/iuq1LQ}
\item Resources: \url{http://goo.gl/51ip9c}
\item List of all the graduate programs and degrees offered: \url{http://goo.gl/3EdsT1}
\end{itemize}

%*********************************************
\addcontentsline{toc}{section}{Congresos y eventos}
\section*{Congresos y eventos}
%*********************************************

\subsection{Workshop on Astrophysics and General Relativity: Astro-GR 2015}
We are happy to announce the 2015 session of the Astro-GR series.
This year we will organize it in Brazil, from August 11 to 15 in Sao Paulo.
I attach more information at the end of this e-mail. Please scroll down.
If you would like to attend, please register as soon as possible here:
\begin{center}
\url{http://ictp-saifr.org/sis/gwavesstep0.php}
\end{center}
and choose "Workshop on Astrophysics and General Relativity: Astro-GR 2015"


%*********************************************
\addcontentsline{toc}{section}{Opinión}
\section*{Opinión}
\subsection{AstroUnicornistis}
Si estás pensando en dedicarte o ya te dedicas a la astronomía o ciencias del espacio ¡piénsalo dos veces! Recientemente se ha encontrado un \textit{virus} conocido como el \textit{AstroUnicornitis} que está atacando a los astrónomos a lo largo del mundo, y al igual que el chikungunya, ya ha llegado al país.\\   

El primer contagiado de este \textit{virus} en nuestro país ha sido el querido doctor Jaime Forero de la U. Andes. El profesor Forero después de tantas horas de fuerte trabajo (como lo demuestra su github) y de pasar toda la noche junto a los `hackeros' organizando el segundo Science Hack Day - Bogotá ha sido visto que ha contraído este espeluznante \textit{virus} como lo demuestra la siguiente imagen 

\begin{center}
\includegraphics[width=.45\textwidth]{AstroUnicornistis.JPG} 
\end{center}

En los próximos volúmenes de RECANEWS estaremos informando sobre la convocatoria que abrirá Colciencias en este pionero tema, en la cual cumplirá con las fechas establecidas (al igual que la de repatriación de cerebros) y donde prometen un gran presupuesto. De igual forma estaremos pendiente del Arxiv.astro-ph por cualquier noticia que pueda interesar a nuestra comunidad. 

%*********************************************





%*****************
\input{notes}
\end{document}
%*****************


